\section{Boolean 相位直线精确界的证明}
\label{bound:appendix}

\subsection{由四面体规则性到带符号分量相等}
Pauli 轨道的局部作用将种子 Bloch 向量 $(x,y,z)$ 变为四个向量
\begin{equation}
 (x,y,z),\quad(-x,-y,z),\quad(x,-y,-z),\quad(-x,y,-z).
 \label{bound:orbit-four}
\end{equation}
局部标签映射覆盖这四种符号模式。不同向量间的内积有三种值：$-x^2-y^2+z^2$、$x^2-y^2-z^2$ 和 $-x^2+y^2-z^2$。它们相等当且仅当 $x^2=y^2=z^2$。所有物理位的半径相同，进一步要求全部 $3n$ 个带符号分量的平方相等。在指定基点，每个带符号分量均等于 $c_n>0$。连续性保证它们在足够小的区间内仍为正，因此分量本身在该区间相等。

对 Boolean 函数 $d$，局部期望中的每一项都带有相位 $e^{it[d(z)-d(z')]}$，其中 $d(z)-d(z')\in\{0,\pm1\}$。因此，任意两个带符号分量之差都形如 $A+B\cos t+C\sin t$。若它在零点邻域恒为零，则在零点处的函数值及前两阶导数给出 $A=B=C=0$。分量相等因而成立于所有实数 $t$，包括可能出现的零半径点。反过来，带符号分量相等就给出式~\eqref{bound:orbit-four} 的规则轨道，半径为 $\sqrt3\,|m_d|$。这证明了定理~\ref{bound:theorem} 使用的从局部到全参数的推论。

\subsection{背景项与系数比较}
记 $N=2^n$、$c=2/N$，并令 $q=d+\delta_0$。对拉回后的带符号 Pauli 算符，写成
\begin{equation}
 R=i^pX^aZ^b,\qquad h=a\oplus b,\qquad
 \Gamma=i^{p-|a|}.
 \label{bound:pullback}
\end{equation}
局部拉回表中的 $a,h$ 均非零，并分为两类：I 类满足 $\Gamma=-i$、$h\cdot a=1$；II 类满足 $\Gamma=-1$、$h\cdot a=0$。利用 $z\ne0$ 时的相位恒等式 $i^{f_0(z)}=-i\,i^{|z|}$，输入种子可写为
\begin{equation}
 -\frac{i}{\sqrt N}\left[
 \sum_z i^{|z|}e^{itq(z)}\ket z
 +(i-e^{it})\ket0\right].
 \label{bound:input-split}
\end{equation}
它对 $R$ 的期望等于背景项
\begin{equation}
 B_R(t)=\frac{\Gamma}{N}\sum_z(-1)^{h\cdot z}
 e^{it[q(z)-q(z\oplus a)]}
 \label{bound:background}
\end{equation}
与交叉项
\begin{equation}
 c\operatorname{Re}\!\left[
 \Gamma e^{-itq(a)}(i-e^{it})\right]
 \label{bound:cross}
\end{equation}
之和。由于 $a\ne0$，零点修正自身没有对角贡献。

在式~\eqref{bound:background} 中将 $z$ 与 $z\oplus a$ 配对，I 类的余弦部分消去。由 $h\ne0$ 知 $\sum_z(-1)^{h\cdot z}=0$，余下项为 $B_I=-W_q(h)\sin t$。对 II 类，配对消去虚部，再利用
\begin{equation}
 \cos\!\bigl(t[q(z)-q(z\oplus a)]\bigr)
 =1-(1-\cos t)D_aq(z)
\end{equation}
得到 $B_{II}=-W_{D_aq}(h)(1-\cos t)/2$。加上式~\eqref{bound:cross}，可得
\begin{equation}
 m_I(t)=\begin{cases}
 c\cos t-W_q(h)\sin t,&q(a)=1,\\
 c-[c+W_q(h)]\sin t,&q(a)=0,
 \end{cases}
 \label{bound:component-i}
\end{equation}
以及
\begin{equation}
 \begin{split}
 m_{II}(t)={}&-\tfrac12W_{D_aq}(h)(1-\cos t)\\
 &+c\bigl\{[1-q(a)]\cos t+q(a)(1-\sin t)\bigr\}.
 \end{split}
 \label{bound:component-ii}
\end{equation}

对 $n\ge3$，两类均非空。式~\eqref{bound:component-i} 中的余弦系数迫使所有 I 类翻转值等于同一个 $s\in\{0,1\}$；式~\eqref{bound:component-ii} 中的正弦系数迫使所有 II 类翻转值等于同一个 $k\in\{0,1\}$。跨类别比较正弦系数得到 $W_q(h)=c(s+k-1)$，比较常数项得到 $W_{D_aq}(h)=2c(s+k-1)$，此时余弦系数也相等。代回即证明式~\eqref{bound:conditions} 的必要性与充分性，并得到四种规律。四者绝对值的最大值给出一般上界。

\subsection{\texorpdfstring{$n=2,3,4$}{n=2,3,4} 时的解析排除}
以下采用 $z_1$ 为最高位的二进制整数编码。对 $n\ge3$ 及每个 I 类字符，Walsh 条件等价于
\begin{equation}
 \sum_z q(z)(-1)^{h\cdot z}=1-s-k,
 \label{bound:linear-character}
\end{equation}
因为 $h\ne0$ 时有 $W_q(h)=-2N^{-1}\sum_zq(z)(-1)^{h\cdot z}$。

当 $n=2$ 时只有 I 类，且其翻转掩码覆盖全部三个非零串。余弦系数相等要求所有 $z\ne0$ 均满足 $q(z)=s$。$s=0$ 与 $s=1$ 分别给出 $q=\delta_0$ 和 $q\equiv1$，共同分量分别为 $c$ 与 $c\cos t$；二者的绝对值均不超过基点值。

当 $n=3$ 时，I 类翻转掩码为 $\{1,2,3,4,5\}$，II 类为 $\{6,7\}$。于是 $q(0)=1$，第一组的函数值均为 $s$，第二组均为 $k$，没有其它未知量。取 $h=1$，式~\eqref{bound:linear-character} 的左端为 $1-s$，得到 $k=0$；取 $h=4$，左端为 $1+s-2k$，得到 $2s-k=0$。故 $s=k=0$，从而 $d=0$。

当 $n=4$ 时，I 类翻转掩码为 $\{1,2,3,8,9\}$，II 类为 $\{6,7,12,13\}$。记
\begin{equation}
 (A,B,C,D,E,F)=(q_4,q_5,q_{10},q_{11},q_{14},q_{15}),
 \label{bound:unknown-values}
\end{equation}
其中 $q_j=q(j)$。四个 I 类字符 $h=1,3,7,15$ 分别给出
\begin{equation}
 \begin{aligned}
 A-B+C-D+E-F+k&=0,\\
 A-B-C+D-E+F+k&=0,\\
 -A+B-C+D+E-F+k&=0,\\
 -A+B+C-D-E+F+k&=0.
 \end{aligned}
 \label{bound:four-boundary-equations}
\end{equation}
四式相加得到 $4k=0$，因此 $k=0$。式~\eqref{bound:four-laws} 只剩 $m=c$ 或 $m=c\cos t$。这一排除涵盖$n=4$ 时的全部 Boolean 函数，未使用任何支持结构假设。它已足以证明半径上界，无需逐一列出可容许支持。

低维上界在 $t=0$ 达到。对每个 $n\ge5$，显式统一生成元满足可容许条件，并给出 $m(-\pi/2)=2c_n$。这两点完成了式~\eqref{bound:sharp-max} 的证明；所得为可达到的最大值，而非仅有上确界。
